% TODO checklist:
% - Inspected: docker-compose.yml (service definitions)
% - Inspected: docker-compose.nginx.yml (production proxy setup)
% - Inspected: Q&A.md sections on deployment assumptions and scaling
% - Inspected: ops/nginx/nginx.conf (upstream definitions)
% - Inspected: src/workers/jobs_worker.py structure
% - Inspected: Dockerfile worker target

\section{Deployment Topologies and Scaling Strategies}
\label{sec:deployment-topologies}

The Cineca Agentic Platform supports multiple deployment topologies ranging from single-node development environments to horizontally-scaled production clusters. This section describes the supported deployment patterns, scaling strategies, and considerations for different operational requirements.

\subsection{Single-Node Deployment}

The simplest deployment topology runs all services on a single host using Docker Compose. This configuration is suitable for:
\begin{itemize}
    \item Development and testing environments
    \item Small-scale production deployments with moderate load
    \item Proof-of-concept installations
    \item Resource-constrained environments (e.g., edge deployments)
\end{itemize}

\begin{verbatim}
┌─────────────────────────────────────────────────────────────────────┐
│                         Single Host                                  │
├─────────────────────────────────────────────────────────────────────┤
│  ┌─────────┐  ┌────────┐  ┌──────────┐  ┌───────┐  ┌─────────────┐ │
│  │ FastAPI │  │ Worker │  │ Postgres │  │ Redis │  │  Memgraph   │ │
│  │   :8000 │  │        │  │   :5432  │  │ :6379 │  │    :7687    │ │
│  └─────────┘  └────────┘  └──────────┘  └───────┘  └─────────────┘ │
│  ┌─────────┐  ┌────────┐  ┌──────────────────────────────────────┐ │
│  │ Ollama  │  │Prometh.│  │         Grafana :3001                │ │
│  │  :11434 │  │  :9090 │  │                                      │ │
│  └─────────┘  └────────┘  └──────────────────────────────────────┘ │
└─────────────────────────────────────────────────────────────────────┘
\end{verbatim}

All services communicate over Docker's internal bridge network (\texttt{app-net}), and data is persisted in named volumes on the host filesystem.

\subsection{Production Deployment with Reverse Proxy}

For production deployments, the platform includes NGINX reverse proxy configuration that provides TLS termination, rate limiting, and load balancing:

\begin{verbatim}
┌───────────────────────────────────────────────────────────────────────┐
│                        Production Deployment                           │
├───────────────────────────────────────────────────────────────────────┤
│                                                                        │
│                    ┌─────────────────────┐                            │
│    Internet ──────▶│  NGINX :443/:80     │                            │
│                    │  TLS Termination    │                            │
│                    │  Rate Limiting      │                            │
│                    └──────────┬──────────┘                            │
│                               │                                        │
│           ┌───────────────────┼───────────────────┐                   │
│           │                   │                   │                   │
│           ▼                   ▼                   ▼                   │
│    ┌─────────────┐    ┌─────────────┐    ┌─────────────────┐         │
│    │  FastAPI    │    │  FastAPI    │    │ Streamlit UI    │         │
│    │  Instance 1 │    │  Instance 2 │    │  Control Panel  │         │
│    └─────────────┘    └─────────────┘    └─────────────────┘         │
│           │                   │                   │                   │
│           └───────────────────┼───────────────────┘                   │
│                               │                                        │
│    ┌──────────────────────────┼──────────────────────────────────┐   │
│    │              Shared Infrastructure                           │   │
│    ├──────────────────────────┼──────────────────────────────────┤   │
│    │  PostgreSQL      Redis        Memgraph       Ollama          │   │
│    └──────────────────────────────────────────────────────────────┘   │
└───────────────────────────────────────────────────────────────────────┘
\end{verbatim}

The NGINX configuration in \texttt{docker-compose.nginx.yml} provides:

\begin{lstlisting}[language=yaml, caption=NGINX production configuration]
services:
  nginx:
    build:
      context: ./ops/nginx
      dockerfile: Dockerfile
    ports:
      - "80:80"
      - "443:443"
    volumes:
      - ./ops/nginx/nginx.conf:/etc/nginx/conf.d/default.conf:ro
      - ./ops/nginx/ssl:/etc/ssl:ro
    depends_on:
      - app
      - ui_control_panel

  app:
    ports: []  # Remove direct port exposure
    environment:
      ENABLE_CORS: "false"       # NGINX handles CORS
      TRUST_PROXY: "true"        # Trust X-Forwarded headers
      SECURE_COOKIES: "true"     # Enable secure cookie flag
\end{lstlisting}

\subsection{Horizontal Scaling Patterns}

\subsubsection{Application Server Scaling}

The FastAPI application is stateless and can be horizontally scaled by running multiple instances behind a load balancer:

\begin{lstlisting}[language=bash, caption=Scaling application instances]
# Scale to 3 application instances
docker compose up -d --scale app=3
\end{lstlisting}

NGINX's upstream configuration supports multiple backend instances:

\begin{lstlisting}[language=nginx, caption=NGINX upstream configuration for load balancing]
upstream app_backend {
    server app:8000 max_fails=3 fail_timeout=30s;
    # Additional instances are automatically discovered
    # by Docker Compose service scaling
}

location /v1/ {
    proxy_pass http://app_backend;
    proxy_http_version 1.1;
    proxy_set_header Host $host;
    proxy_set_header X-Real-IP $remote_addr;
    proxy_set_header X-Forwarded-For $proxy_add_x_forwarded_for;
    proxy_set_header X-Forwarded-Proto $scheme;
}
\end{lstlisting}

\subsubsection{Worker Process Scaling}

Background job workers can be scaled independently from the API servers:

\begin{lstlisting}[language=bash, caption=Scaling worker processes]
# Scale to 5 worker instances
docker compose up -d --scale worker=5
\end{lstlisting}

Workers implement distributed locking through Redis to prevent duplicate job processing:
\begin{itemize}
    \item Jobs are claimed atomically using \texttt{SETNX} operations
    \item Heartbeats indicate worker liveness
    \item Orphaned jobs are recovered after missed heartbeats
\end{itemize}

\subsubsection{Redis Scaling Considerations}

For high-traffic deployments, Redis can be scaled using:

\begin{enumerate}
    \item \textbf{Redis Sentinel}: Provides automatic failover for high availability
    \item \textbf{Redis Cluster}: Provides horizontal sharding for large datasets
    \item \textbf{Managed Redis}: Cloud-managed solutions (e.g., AWS ElastiCache, Azure Cache for Redis)
\end{enumerate}

The platform's Redis usage is compatible with Redis Cluster mode, with the following considerations:
\begin{itemize}
    \item Job queues use \texttt{BRPOP} which requires single-key operations
    \item Rate limiting uses sorted sets with tenant-scoped keys
    \item Session state is stored with TTL-based expiration
\end{itemize}

\subsection{Kubernetes Deployment Patterns}

While the primary deployment method uses Docker Compose, the platform is designed to be Kubernetes-ready:

\paragraph{Health Probe Compatibility}
The platform exposes Kubernetes-compatible health endpoints:

\begin{lstlisting}[language=yaml, caption=Kubernetes probe configuration example]
livenessProbe:
  httpGet:
    path: /v1/health/live
    port: 8000
  initialDelaySeconds: 30
  periodSeconds: 10

readinessProbe:
  httpGet:
    path: /v1/health/ready
    port: 8000
  initialDelaySeconds: 10
  periodSeconds: 5

startupProbe:
  httpGet:
    path: /v1/health/startup
    port: 8000
  failureThreshold: 30
  periodSeconds: 10
\end{lstlisting}

\paragraph{ConfigMap and Secret Management}
Configuration can be injected via Kubernetes ConfigMaps and Secrets:

\begin{lstlisting}[language=yaml, caption=Kubernetes environment configuration example]
env:
  - name: DB_HOST
    valueFrom:
      configMapKeyRef:
        name: cineca-config
        key: db-host
  - name: DB_PASSWORD
    valueFrom:
      secretKeyRef:
        name: cineca-secrets
        key: db-password
  - name: OIDC_ISSUER
    valueFrom:
      configMapKeyRef:
        name: cineca-config
        key: oidc-issuer
\end{lstlisting}

\paragraph{Horizontal Pod Autoscaling}
The FastAPI application exports Prometheus metrics that can drive Kubernetes HPA decisions:

\begin{lstlisting}[language=yaml, caption=Horizontal Pod Autoscaler example]
apiVersion: autoscaling/v2
kind: HorizontalPodAutoscaler
metadata:
  name: cineca-app-hpa
spec:
  scaleTargetRef:
    apiVersion: apps/v1
    kind: Deployment
    name: cineca-app
  minReplicas: 2
  maxReplicas: 10
  metrics:
    - type: Resource
      resource:
        name: cpu
        target:
          type: Utilization
          averageUtilization: 70
\end{lstlisting}

\subsection{Scaling Strategy Summary}

Table~\ref{tab:scaling-strategies} summarizes the scaling approaches for each component.

\begin{table}[ht]
\centering
\caption{Scaling strategies by component.}
\label{tab:scaling-strategies}
\small
\begin{tabular}{lp{4cm}p{5cm}}
\hline
\textbf{Component} & \textbf{Scaling Approach} & \textbf{Considerations} \\
\hline
FastAPI App & Horizontal (replicas) & Stateless; requires shared Redis for sessions \\
Worker & Horizontal (replicas) & Distributed locking prevents duplicate processing \\
PostgreSQL & Vertical or read replicas & Consider managed solutions for HA \\
Redis & Sentinel or Cluster & Job queues require single-key operations \\
Memgraph & Vertical & Graph analytics benefit from memory \\
Ollama & Vertical (GPU/memory) & Model loading is memory-intensive \\
NGINX & Horizontal (replicas) & Typically not the bottleneck \\
\hline
\end{tabular}
\end{table}

\subsection{Resource Requirements}

Table~\ref{tab:resource-requirements} provides recommended resource allocations for different deployment scales.

\begin{table}[ht]
\centering
\caption{Recommended resource allocations by deployment scale.}
\label{tab:resource-requirements}
\small
\begin{tabular}{lcccc}
\hline
\textbf{Component} & \textbf{Dev} & \textbf{Small Prod} & \textbf{Medium Prod} & \textbf{Large Prod} \\
\hline
FastAPI App & 512MB / 0.5 CPU & 1GB / 1 CPU & 2GB / 2 CPU & 4GB / 4 CPU \\
Worker & 256MB / 0.25 CPU & 512MB / 0.5 CPU & 1GB / 1 CPU & 2GB / 2 CPU \\
PostgreSQL & 256MB / 0.5 CPU & 1GB / 1 CPU & 4GB / 2 CPU & 8GB / 4 CPU \\
Redis & 128MB / 0.25 CPU & 256MB / 0.5 CPU & 1GB / 1 CPU & 4GB / 2 CPU \\
Memgraph & 512MB / 0.5 CPU & 2GB / 1 CPU & 8GB / 2 CPU & 32GB / 4 CPU \\
Ollama (CPU) & 7GB / 4 CPU & 10GB / 8 CPU & --- & --- \\
Ollama (GPU) & --- & 8GB VRAM & 16GB VRAM & 24GB+ VRAM \\
\hline
\end{tabular}
\end{table}

\subsection{High Availability Considerations}

For high-availability deployments, consider:

\begin{enumerate}
    \item \textbf{Multiple API instances}: At least 2 instances behind a load balancer
    \item \textbf{Database replication}: PostgreSQL streaming replication or managed database service
    \item \textbf{Redis redundancy}: Redis Sentinel for automatic failover
    \item \textbf{Health-based routing}: Remove unhealthy instances from rotation
    \item \textbf{Geographic distribution}: Deploy across availability zones
    \item \textbf{Backup strategy}: Regular automated backups with tested restore procedures
\end{enumerate}

The platform's graceful degradation features (circuit breakers, fallback modes) help maintain partial functionality during component failures, enhancing overall system resilience.

